\subsection{B.2. Arbitrary algebraic residue extensions}

We use an algebraic extension in the literal sense: $k\subset K$ are fields,
with no bound on the cardinality of $K$, and every element of $K$ is algebraic
over $k$.  The following resolves the question on cote 23, scan 063.  The
word ``every'' in its conclusion includes lifts that introduce new variables
with coefficients in the maximal ideal; retaining triangularity upstairs is
not required.

\begin{lemma}[Proved: lifting a finite regular sequence]\label{SA23-B2-1}
Let $(A,\mathfrak m,k)$ be an Artinian local ring, let
$P=A[T_\alpha\mid\alpha\in\Lambda]$, and let $f_1,\ldots,f_s\in P$.
If their images in $\overline P=k[T_\alpha]$ form a regular sequence, then
$P/(f_1,\ldots,f_s)$ is free over $A$.  The $f_i$ also form a
Koszul-regular sequence in $P$.
\end{lemma}
\begin{proof}
Put $D_\bullet=K_\bullet(P;f_1,\ldots,f_s)$, with the Koszul differential
in homological degree $-1$.  Filter each term by $\mathfrak m^jD_q$.
Because $P$ is $A$-free, its associated graded complex in filtration degree
$j$ is
\[
 (\mathfrak m^j/\mathfrak m^{j+1})\otimes_k
 K_\bullet(\overline P;\overline f_1,\ldots,\overline f_s).
\]
The latter has homology only in degree zero.  The filtration is finite:
$\mathfrak m^h=0$ for some integer $h$.  A direct successive lifting of
cycles, starting at the lowest nonzero filtration degree, therefore proves
$H_q(D)=0$ for $q>0$.  The same argument for boundaries gives
\[
 \operatorname{gr}_{\mathfrak m}\bigl(P/(f_1,\ldots,f_s)\bigr)
 \simeq \operatorname{gr}_{\mathfrak m}(A)
       \otimes_k \overline P/(\overline f_1,\ldots,\overline f_s).
\]
More explicitly, a cycle whose leading term is a boundary can be changed
by a boundary to increase its filtration; after at most $h$ changes it is
zero.  Applying this to a degree-zero element that becomes zero in the
quotient proves the asserted equality of associated graded modules.

Choose a $k$-basis $(b_\gamma)_{\gamma\in\Gamma}$ of the residue quotient
and lift it to elements $\widetilde b_\gamma$ of $P/(f_1,\ldots,f_s)$.
The $A$-linear map $A^{(\Gamma)}\longrightarrow P/(f_1,\ldots,f_s)$
sending the standard basis to these lifts induces an isomorphism on the
associated graded modules just displayed.  A map with a finite filtration
that induces an isomorphism on every associated graded piece is an
isomorphism: injectivity and surjectivity both follow by raising the
filtration at most $h$ times.  This argument works for an infinite basis,
since every vector and every correction is a finite linear combination.
\end{proof}

\begin{theorem}[Proved: Deligne's algebraic-extension question]\label{SA23-B2-2}
For every algebraic extension $K/k$ there are sets of variables and
relations such that
\[
 K\simeq k[T_\alpha\mid\alpha<\lambda]/(P_\alpha\mid\alpha<\lambda),
\]
and for every Artinian local ring $(A,\mathfrak m)$ supplied with an
identification $A/\mathfrak m=k$, and every family of polynomials
$\widetilde P_\alpha\in A[T_\beta\mid\beta<\lambda]$ reducing to
$P_\alpha$, the algebra
\[
 A[T_\beta\mid\beta<\lambda]/(\widetilde P_\alpha\mid\alpha<\lambda)
\]
is flat over $A$.  It is, in fact, free as an $A$-module.
\end{theorem}
\begin{proof}
Well-order a set of generators of $K$ over $k$, say
$(x_\alpha)_{\alpha<\lambda}$.  Write
$K_{<\alpha}=k[x_\beta\mid\beta<\alpha]$.  This ring is a field:
every element belongs to a finite algebraic extension of $k$ contained in
it, so every nonzero element has its inverse in the same finite extension.
Let $Q_\alpha(X)\in K_{<\alpha}[X]$ be the monic minimal polynomial of
$x_\alpha$.  Represent each of its finitely many coefficients by a
polynomial in finitely many preceding generators.  The resulting relation
$P_\alpha$ is monic in $T_\alpha$ and involves no variable $T_\beta$
with $\beta>\alpha$.

Transfinite induction shows that the quotient by all relations with index
less than $\alpha$ is $K_{<\alpha}$: at a successor stage use the minimal
polynomial, and at a limit stage take the union.  This proves the asserted
presentation of $K$.

Any finite subset $\alpha_1<\cdots<\alpha_s$ of the relations is a
regular sequence in $k[T_\beta\mid\beta<\lambda]$.  Before adjoining
$P_{\alpha_i}$, the earlier relations do not contain $T_{\alpha_i}$;
the ambient quotient is consequently still a polynomial ring in
$T_{\alpha_i}$ over a ring.  A monic polynomial in this variable is a
nonzerodivisor over every coefficient ring, by comparing the leading
coefficient of a product.  The quotient is nonzero, for example because
it maps to $K$.  Lemma~\ref{SA23-B2-1} now says that the quotient of
$A[T_\beta]$ by the corresponding finite set of arbitrary lifted
relations is $A$-free.  The quotient by all relations is their filtered
colimit.  Tensor product commutes with filtered colimits and filtered
colimits of modules preserve exact sequences, so this quotient is flat.

For completeness, a flat module $M$ over an Artinian local ring $A$ is
free.  Lift a basis of $M/\mathfrak mM$ to obtain
$F=A^{(\Gamma)}\to M$.  Its cokernel equals its own product with
$\mathfrak m$ and hence is zero, since $\mathfrak m$ is nilpotent.
If $L$ is the kernel, flatness gives an injection
$L/\mathfrak mL\to F/\mathfrak mF$; this map is zero because
$F/\mathfrak mF\to M/\mathfrak mM$ is an isomorphism.  Thus
$L=\mathfrak mL=0$.  No finite-generation hypothesis was used.
\end{proof}

\begin{remark}[Scope and literature search]\label{SA23-B2-3}
The construction does not require a separating transcendence basis,
a finite $p$-basis, a finite exponent, or a simple extension.  It settles
the algebraic question stated on scan 063, not a different assertion
about all ring maps.  Our search for algebraic field extensions with
regular-sequence presentations did not identify a priority claim for
this elementary transfinite construction; we make none.  The flat
lifting mechanism is the familiar fiberwise regular-sequence criterion;
we supplied a proof valid for the infinitely generated polynomial rings
needed here.
\end{remark}
